\section{Methodology}
\label{sec:method}

In this section, we present the formal mathematical framework of \textbf{ChebyMerge}. First, we establish the problem setup of multi-task model merging under test-time adaptation. Second, we introduce the spectral parameterization using Chebyshev polynomials of the first kind. Third, we formulate the unsupervised test-time optimization objective. Finally, we provide a rigorous theoretical proof demonstrating the superior numerical conditioning of ChebyMerge over monomial-based continuous alternatives like PolyMerge.

\subsection{Problem Formulation and Notation}
Let $L$ denote the number of layers (or parameter blocks) in a deep neural network, and let $K$ denote the number of task-specific expert models. We assume all expert models share a common pre-trained base model architecture, parameterized by $\Theta_{\text{base}}$. Each task expert $k \in \{1, \dots, K\}$ is parameterized by $\Theta_k$, representing weights fine-tuned on task $k$ starting from $\Theta_{\text{base}}$. Following the task arithmetic framework \cite{ilharco2022editing}, we construct the task-specific direction vectors---termed \emph{task vectors}---for each layer $l \in \{0, 1, \dots, L-1\}$ as:
\begin{equation}
\mathbf{\Delta}_{k, l} = \Theta_{k, l} - \Theta_{\text{base}, l}
\end{equation}
where $\Theta_{k, l}$ and $\Theta_{\text{base}, l}$ represent the parameter tensors of task $k$ and the base model at layer $l$, respectively. 

Rather than applying a uniform merging coefficient across all layers, we parameterize the merging coefficients layer-wise. Let $\lambda_{k, l}$ denote the merging coefficient applied to task vector $k$ at layer $l$. The consolidated weights for layer $l$, denoted by $\Theta_{\text{merged}, l} \in \mathbb{R}^{M_l}$, are defined as:
\begin{equation}
\Theta_{\text{merged}, l}(\boldsymbol{\lambda}_l) = \Theta_{\text{base}, l} + \sum_{k=1}^K \lambda_{k, l} \mathbf{\Delta}_{k, l}
\end{equation}
where $\boldsymbol{\lambda}_l = [\lambda_{1, l}, \dots, \lambda_{K, l}]^T \in \mathbb{R}^K$. Under unconstrained dynamic merging (e.g., AdaMerging \cite{lu2023adamerging}), the optimizer treats each $\lambda_{k, l}$ as an independent free variable, resulting in $K \times L$ parameters to optimize.

\subsection{Continuous Subspace Projection via Chebyshev Polynomials}
To prevent transductive overfitting on small test batches, we project the high-dimensional coefficient space onto a low-dimensional continuous subspace. We model the layer-wise coefficients as continuous functions of network depth, parameterized via orthogonal polynomials. 

\subsubsection{Linear Coordinate Mapping}
We map the discrete layer indices $l \in \{0, 1, \dots, L-1\}$ to the compact Chebyshev domain $x \in [-1, 1]$ via the linear transformation:
\begin{equation}
x_l = \frac{2l}{L-1} - 1, \quad \forall l \in \{0, 1, \dots, L-1\}
\end{equation}

\subsubsection{Chebyshev Recurrence Relation}
We define the $j$-th Chebyshev polynomial of the first kind, $T_j(x)$, recursively as:
\begin{align}
T_0(x) &= 1 \\
T_1(x) &= x \\
T_j(x) &= 2x T_{j-1}(x) - T_{j-2}(x), \quad \forall j \ge 2
\end{align}
Equivalently, Chebyshev polynomials can be defined trigonometrically as $T_j(x) = \cos(j \arccos(x))$ for $x \in [-1, 1]$.

\subsubsection{Subspace Parameterization}
We parameterize the spatial merging coefficient $\lambda_{k, l}$ for task $k$ and layer $l$ as a linear combination of Chebyshev basis functions evaluated at $x_l$:
\begin{equation}
\label{eq:cheby_param}
\lambda_{k, l}(\boldsymbol{\alpha}) = \sum_{j=0}^d \alpha_{k, j} T_j(x_l)
\end{equation}
where $\boldsymbol{\alpha} = \{ \alpha_{k, j} \} \in \mathbb{R}^{K \times (d+1)}$ represents the matrix of learnable spectral parameters, and $d$ is the polynomial degree ($d \ll L$, typically $d \in \{1, 2, 3\}$).

This polynomial representation has an elegant spectral interpretation rooted in digital signal processing. Under the coordinate transformation $x = \cos(\theta)$, the continuous Chebyshev polynomial of the first kind is isomorphic to the continuous Cosine Transform. On a discrete grid, when using cosine-spaced nodes (Chebyshev-Gauss-Lobatto nodes), Chebyshev expansion is exactly isomorphic to the Discrete Cosine Transform (DCT) of Type I:
\begin{equation}
T_j(\cos(\theta)) = \cos(j \theta)
\end{equation}
In our case, because we evaluate the basis functions on a \emph{uniform discrete grid} $x_l \in [-1, 1]$, this exact discrete isomorphism is lost, and we instead obtain a warped spectral representation. Evaluating Chebyshev polynomials on a uniform grid introduces a non-linear coordinate mapping $\theta_l = \arccos(x_l) = \arccos\left(\frac{2l}{L-1} - 1\right)$, which warps the local spatial frequencies of the basis functions. Because the derivative of $\arccos(x)$ is $- (1 - x^2)^{-1/2}$, the local spatial frequency is stretched near the boundaries ($x \approx \pm 1$) and compressed in the middle ($x \approx 0$). Far from being a limitation, this "frequency warping" acts as a highly beneficial, \emph{foveated spectral filter} that perfectly aligns with the physical sensitivity profile of deep models. Near the early and late boundary layers (where sensitivities are high and exhibit rapid spatial variations), the warped basis provides fine-grained spatial resolution, allowing the model to capture high-frequency components. In contrast, in the intermediate layers (where sensitivities are flat and robust), the warped basis applies an aggressive low-pass filter, suppressing high-frequency noise and preventing the optimizer from fitting transductive local fluctuations. This foveated spectral filtering represents a compelling synthesis of numerical analysis and physical network prior. By eliminating these high-frequency components in intermediate layers while preserving local details at the boundaries, we successfully filter out white transductive noise from the test-time adaptation stream, while keeping the underlying layer sensitivity profile intact.

We construct a precomputed Chebyshev design matrix $\mathbf{C} \in \mathbb{R}^{L \times (d+1)}$, where:
\begin{equation}
C_{l, j} = T_j(x_l)
\end{equation}
The mapping from the spectral parameters to the spatial coefficients can be compactly written using matrix multiplication:
\begin{equation}
\boldsymbol{\lambda}_k = \mathbf{C} \boldsymbol{\alpha}_k^T \in \mathbb{R}^L
\end{equation}
where $\boldsymbol{\alpha}_k = [\alpha_{k, 0}, \dots, \alpha_{k, d}] \in \mathbb{R}^{1 \times (d+1)}$ is the spectral parameter vector for task $k$. The precomputed design matrix $\mathbf{C}$ is registered as a constant PyTorch buffer during model initialization, ensuring that the projection from spectral to spatial domains introduces zero computational overhead during optimization.

\subsection{Unsupervised Test-Time Adaptation Objective}
During test-time adaptation, the model receives unlabeled local data batches from task-specific streams $\mathcal{D}_k^{\text{unlabeled}}$. We optimize the spectral parameters $\boldsymbol{\alpha}$ on-the-fly by minimizing the Shannon entropy of the merged model's predictions. 

Let $f_{\Theta_{\text{merged}}(\boldsymbol{\alpha})}(x) \in \mathbb{R}^C$ denote the softmax class probability output of the merged model for input $x$, where $C$ is the number of classes. The self-supervised test-time adaptation loss $\mathcal{L}_{\text{TTA}}(\boldsymbol{\alpha})$ is defined as:
\begin{equation}
\mathcal{L}_{\text{TTA}}(\boldsymbol{\alpha}) = \sum_{k=1}^K \mathbb{E}_{x \sim \mathcal{D}_k^{\text{unlabeled}}} \left[ H\left(f_{\Theta_{\text{merged}}(\boldsymbol{\alpha})}(x)\right) \right]
\end{equation}
where the Shannon entropy $H(\mathbf{p})$ is evaluated as:
\begin{equation}
H(\mathbf{p}) = -\sum_{c=1}^C p_c \log (p_c + \epsilon)
\end{equation}
and $\epsilon = 10^{-8}$ is a small numerical stabilizer. 

\subsection{Theoretical Analysis of Numerical Conditioning}
We now provide a rigorous mathematical justification for the superior stability of ChebyMerge compared to monomial-based continuous parameterizations (such as PolyMerge).

Let $\mathbf{X} \in \mathbb{R}^{L \times (d+1)}$ represent the design matrix that maps spectral parameters to spatial coefficients. Under PolyMerge, the spatial coefficients are modeled as:
\begin{equation}
\lambda_{k, l}(\boldsymbol{\gamma}) = \sum_{j=0}^d \gamma_{k, j} \bar{l}^j
\end{equation}
where $\bar{l} = \frac{l}{L-1} \in [0, 1]$ is the normalized layer depth. The mapping is mediated by the standard Vandermonde-type monomial design matrix $\mathbf{V} \in \mathbb{R}^{L \times (d+1)}$, where $V_{l, j} = \bar{l}^j$.

During gradient descent, the local curvature and conditioning of the loss landscape with respect to the learnable spectral parameters is characterized by the local Hessian matrix. By applying the chain rule, the gradient of the loss with respect to the spectral parameter vector $\boldsymbol{\alpha}_k$ is:
\begin{equation}
\nabla_{\boldsymbol{\alpha}_k} \mathcal{L} = \mathbf{X}^T \nabla_{\boldsymbol{\lambda}_k} \mathcal{L}
\end{equation}
Consequently, the Hessian of the loss with respect to the spectral parameters $\mathbf{H}_{\boldsymbol{\alpha}_k}$ is related to the spatial Hessian $\mathbf{H}_{\boldsymbol{\lambda}_k}$ via:
\begin{equation}
\mathbf{H}_{\boldsymbol{\alpha}_k} = \mathbf{X}^T \mathbf{H}_{\boldsymbol{\lambda}_k} \mathbf{X}
\end{equation}
Assuming the spatial Hessian $\mathbf{H}_{\boldsymbol{\lambda}_k}$ is locally well-behaved (e.g., bounded), the condition number of the optimization metric $\mathbf{H}_{\boldsymbol{\alpha}_k}$ is directly controlled by the condition number of the Gram matrix:
\begin{equation}
\kappa(\mathbf{X}^T \mathbf{X}) = \frac{\sigma_{\max}(\mathbf{X}^T \mathbf{X})}{\sigma_{\min}(\mathbf{X}^T \mathbf{X})}
\end{equation}
where $\sigma_{\max}$ and $\sigma_{\min}$ represent the maximum and minimum singular values, respectively.

\begin{theorem}[Condition Number of Monomial Basis]
The condition number of the monomial Gram matrix $\mathbf{V}^T \mathbf{V}$ grows exponentially with the polynomial degree $d$:
\begin{equation}
\kappa(\mathbf{V}^T \mathbf{V}) = \mathcal{O}(4^d)
\end{equation}
\end{theorem}

\begin{proof}
In the continuous limit as $L \to \infty$, the discrete elements of the monomial Gram matrix $\mathbf{G}_V = \mathbf{V}^T \mathbf{V}$ scale proportional to the inner products of the monomial basis functions $\{1, x, x^2, \dots, x^d\}$ over the domain $x \in [0, 1]$:
\begin{equation}
(G_V)_{i, j} \propto \int_0^1 x^i x^j \, dx = \int_0^1 x^{i+j} \, dx = \frac{1}{i+j+1}
\end{equation}
The continuous limit of the monomial Gram matrix is the celebrated \emph{Hilbert matrix} $\mathbf{H}_d$, whose elements are $(H_d)_{i, j} = \frac{1}{i+j-1}$ (for 1-based indexing). It is a well-established result in numerical analysis that the condition number of the Hilbert matrix grows exponentially with its dimension \cite{gautschi1990influence}:
\begin{equation}
\kappa(\mathbf{H}_d) \sim \mathcal{O}\left( e^{3.525 d} \right) \approx \mathcal{O}(4^d)
\end{equation}
This extreme collinearity arises because as the degree $d$ increases, the monomial basis functions $x^j$ become highly indistinguishable over $[0, 1]$, resulting in a nearly singular Gram matrix.
\end{proof}

\begin{theorem}[Condition Number of Chebyshev Basis]
The condition number of the Chebyshev Gram matrix $\mathbf{C}^T \mathbf{C}$ is bounded by a small constant close to 1:
\begin{equation}
\kappa(\mathbf{C}^T \mathbf{C}) \le B \ll \kappa(\mathbf{V}^T \mathbf{V})
\end{equation}
where $B$ is independent of the number of layers $L$ for $d \ll L$.
\end{theorem}

\begin{proof}
Chebyshev polynomials of the first kind $T_j(x)$ are strictly orthogonal with respect to the continuous weighted inner product on $x \in [-1, 1]$ under the Chebyshev weight function $w(x) = (1 - x^2)^{-1/2}$:
\begin{equation}
\int_{-1}^1 T_i(x) T_j(x) \frac{1}{\sqrt{1-x^2}} \, dx = \begin{cases} 0 & \text{if } i \ne j \\ \pi & \text{if } i = j = 0 \\ \pi / 2 & \text{if } i = j > 0 \end{cases}
\end{equation}
When evaluated on a discrete grid $x_l \in [-1, 1]$, the Chebyshev polynomials preserve their discrete orthogonality exactly if evaluated over Chebyshev-Gauss-Lobatto (cosine-spaced) nodes. For a uniform discrete grid $x_l = \frac{2l}{L-1} - 1$, this exact discrete orthogonality is lost, and the Gram matrix is only \emph{approximately} diagonal. Nonetheless, because the basis functions $T_j(x)$ oscillate orthogonally across the interval $[-1, 1]$, the off-diagonal elements $(G_C)_{i, j}$ for $i \ne j$ are exceptionally small, while the diagonal elements remain well-scaled. Consequently, although the uniform discrete grid spacing invalidates strict mathematical diagonal dominance for arbitrary dimensions and polynomial degrees, the off-diagonal entries remain exceptionally small and tightly bounded relative to the diagonal. Rather than relying on strict mathematical diagonal dominance, we emphasize that the eigenvalues of the Chebyshev Gram matrix $\mathbf{C}^T \mathbf{C}$ remain extremely tightly clustered near their continuous limits. This bounds the condition number:
\begin{equation}
\kappa(\mathbf{C}^T \mathbf{C}) \approx \frac{\pi / 2}{\pi / 2} \approx 1
\end{equation}
yielding an exceptionally isotropic and well-conditioned optimization landscape.
\end{proof}

This conditioning breakthrough has a profound impact on optimization. In PolyMerge, the massive eigenvalue spread forces gradient updates to either vanish along high-frequency directions or oscillate violently along low-frequency directions, requiring highly specialized learning rate tuning. ChebyMerge's near-perfect conditioning decouples the gradient updates for each spectral coefficient, enabling isotropic, stable, and rapid gradient flow.

\subsection{Implicit Boundary Sensitivity Matching}
Beyond numerical stability, ChebyMerge incorporates a physical structural prior that aligns with deep network properties. It is widely known that the sensitivity of deep models is highly non-uniform across layers: early layers (responsible for feature extraction) and deep layers (responsible for projection and classification) are highly sensitive to weight perturbations, whereas intermediate layers exhibit a high degree of flatness and robustness \cite{zhang2019all}.

The roots and extrema of the Chebyshev polynomials $T_j(x) = \cos(j \arccos(x))$ are known to cluster densely near the boundaries of the interval $[-1, 1]$. By mapping the network depth to this interval, the Chebyshev basis naturally allocates higher spatial sampling density and representation resolution to the early ($l \approx 0 \implies x \approx -1$) and deep ($l \approx L-1 \implies x \approx 1$) layers. This implicit concentration of resolution matches the network's physical sensitivity profile perfectly, allowing ChebyMerge to represent complex, sharp sensitivity transitions at the boundaries using a very low polynomial degree ($d=2$), where standard monomials would suffer from boundary oscillations (Runge's phenomenon).

\subsection{The Conditioning-Generalization Paradox and the Principle of Controllable Regularization}
\label{sec:paradox}

We now address a profound and highly subtle conceptual phenomenon arising from the intersection of numerical analysis and statistical learning theory in on-the-fly optimization, which we term the \textbf{Conditioning-Generalization Paradox}. 

Recall that during test-time adaptation, the optimization objective is heavily corrupted by transductive local sampling noise $\boldsymbol{\eta}$. In PolyMerge, the spatial coefficients are parameterized via a monomial basis $\mathbf{V}$, which we proved possesses an exponentially large condition number:
\begin{equation}
\kappa(\mathbf{V}^T \mathbf{V}) = \mathcal{O}(4^d)
\end{equation}
From an optimization perspective, this extreme ill-conditioning implies that the eigenvalues of the Gram matrix span multiple orders of magnitude. The singular values corresponding to higher-order terms (e.g., $\bar{l}^2, \bar{l}^3$) are extremely close to zero, effectively compressing the local gradient fields along those directions. Consequently, during gradient descent with a standard learning rate, the updates to the higher-order spectral coefficients are heavily damped:
\begin{equation}
\mathbf{H}_{\boldsymbol{\gamma}_k} = \mathbf{V}^T \mathbf{H}_{\boldsymbol{\lambda}_k} \mathbf{V}
\end{equation}
This behaves as an uncontrolled, implicit low-pass spectral filter---often termed \emph{implicit spectral damping} or \emph{implicit early stopping}---that locks the higher-degree components of the polynomial and prevents them from adapting to high-frequency transductive noise.

A subtle but crucial question arises: since the Adam optimizer employs coordinate-wise learning rate scaling (dividing the first moment by the root-mean-square of second moments), should it not neutralize this ill-conditioning by scaling up updates in the "stiff" directions? Theoretically, yes, but in practice, the presence of transductive noise prevents this. Along the highly ill-conditioned (stiff) directions, the true gradient signal has an extremely low Signal-to-Noise Ratio (SNR). When computing Adam's update $\Delta \gamma_j \propto \hat{m}_j / (\sqrt{\hat{v}_j} + \epsilon)$, the denominator $\hat{v}_j$ (which tracks the squared gradients) is dominated by the variance of the transductive noise $\boldsymbol{\eta}$, while the first moment $\hat{m}_j$ (which averages out the zero-mean noise) remains close to zero. Consequently, the noise-dominated denominator heavily suppresses the update in these directions, preventing Adam from scaling up the true signal. Thus, the implicit spectral damping of the ill-conditioned monomial basis persists even under adaptive optimizers like Adam, acting as an accidental, noise-driven regularizer.

In contrast, because ChebyMerge's design matrix $\mathbf{C}$ is exceptionally well-conditioned ($\kappa(\mathbf{C}^T \mathbf{C}) \approx 2.95$ for cubic), the optimization landscape is isotropic. This means that all spectral coefficients---including higher-order terms representing more rapid spatial variations---can be updated with equal ease and speed. While this isotropic gradient flow enables rapid and stable convergence, it also grants the optimizer the freedom to fit high-frequency transductive noise on noisy local batches, leading to a slight overfit compared to the heavily damped monomial basis at higher polynomial degrees.

We argue, however, that relying on extreme numerical ill-conditioning and matrix singularity as an implicit regularizer is sub-optimal, unpredictable, and fragile for three major reasons:
\begin{enumerate}
    \item \textbf{Extreme Sensitivity to Hyperparameters:} Implicit spectral damping is highly sensitive to the exact choice of learning rate, optimization step count, and optimizer initialization. A minor change in these parameters can cause the damped coefficients to suddenly explode or fail to adapt at all, rendering the optimization highly unpredictable.
    \item \textbf{Gradient Vanishing and Loss of Expressivity:} For models with larger depth $L$ or under fine-grained block-wise merging, the ill-conditioning of monomials worsens to the point where higher-order coefficients suffer from numerical underflow, destroying the model's capacity to represent any complex sensitivity profiles.
    \item \textbf{Loss of Controllability:} Relying on accidental numerical errors conflates optimization stability with parameter regularization. 
\end{enumerate}

To achieve scientific rigor, a method must separate numerical conditioning from regularization. ChebyMerge adheres to the \textbf{Principle of Controllable Regularization}: by maintaining perfect numerical conditioning, we guarantee a stable and isotropic optimization metric, and then enforce regularization explicitly and controllably. This can be achieved through:
\begin{itemize}
    \item \textbf{Spectral Dimension Capping (Degree Capping):} Keeping the polynomial degree strictly bounded ($d=1$ or $d=2$).
    \item \textbf{Explicit Spectral Weight Decay (L2 Penalization):} Adding a controlled ridge penalty on the higher-order spectral parameters:
    \begin{equation}
    \mathcal{L}_{\text{reg}}(\boldsymbol{\alpha}) = \gamma_{\text{spectral}} \sum_{j=1}^d \alpha_{k, j}^2
    \end{equation}
    where $\gamma_{\text{spectral}}$ explicitly controls the damping of high-frequency components without compromising the numerical conditioning of the landscape.
    \item \textbf{Controllable Spectral Decay (CSD):} Instead of applying a uniform learning rate across all spectral components, we can explicitly scale the learning rate of each Chebyshev coefficient proportional to its frequency order. Specifically, for the $j$-th spectral coefficient $\alpha_{k, j}$, we apply a decayed learning rate:
    \begin{equation}
    \eta_j = \eta_{\text{base}} \cdot \gamma_{\text{CSD}}^j
    \end{equation}
    where $\gamma_{\text{CSD}} \in (0, 1]$ is a controllable spectral decay factor (e.g., $\gamma_{\text{CSD}} = 0.2$). This mimics the damping of high-frequency terms achieved implicitly by PolyMerge, but does so in a completely explicit, controllable, and predictable manner. The underlying optimization landscape remains perfectly well-conditioned and isotropic, guaranteeing that we maintain perfect coordinate stability while selectively filtering out transductive noise in higher frequencies.
\end{itemize}
This principled decoupling of optimization stability and parameter regularization establishes ChebyMerge as a mathematically superior and far more reliable framework for dynamic model merging.
